\documentclass[12pt]{article}

% --- Idioma ---
\usepackage[spanish,es-nodecimaldot]{babel}

% --- Fuente Times New Roman (compilar con XeLaTeX) ---
\usepackage{fontspec}
\setmainfont{Times New Roman}

% --- Márgenes 2.5 cm por lado ---
\usepackage[top=2.5cm, bottom=2.5cm, left=2.5cm, right=2.5cm]{geometry}

% --- Interlineado 1.5 ---
\usepackage{setspace}
\onehalfspacing

% --- Espaciado posterior 6 pt, sin sangría de párrafo ---
\setlength{\parindent}{0pt}
\setlength{\parskip}{6pt}

% --- Sangría francesa en referencias (APA) ---
\usepackage{hanging}

% --- Hipervínculos ---
\usepackage[hidelinks]{hyperref}

% ============================================================
\begin{document}
% ============================================================

\section{Marco Teórico}

\subsection{La salud como capacidad fundamental y dimensión de la pobreza multidimensional}

El análisis económico del sector salud requiere partir de una distinción conceptual previa: la salud
no es únicamente un bien de consumo ni un insumo productivo, sino una capacidad fundamental del
ser humano. Esta distinción tiene implicancias directas sobre cómo se mide el bienestar y sobre qué
papel le corresponde al Estado en la provisión de servicios sanitarios.

Sen (1988) argumenta que la enfermedad no es simplemente un síntoma observable de pobreza,
sino una restricción a las capacidades básicas de las personas: priva al individuo de su libertad para
funcionar, trabajar y participar en la vida social. Este argumento es ampliado en Sen (2000), donde
el desarrollo económico es reinterpretado no como crecimiento del producto per cápita, sino como
expansión de las libertades reales que los individuos pueden ejercer. Bajo este enfoque, un país que
duplica su PIB sin reducir la mortalidad infantil o aumentar la esperanza de vida no puede
considerarse plenamente desarrollado. La salud, junto con la educación y las libertades políticas,
constituye una de las capacidades centrales cuya expansión define el progreso humano genuino.

Esta perspectiva tiene consecuencias metodológicas importantes. Si el bienestar no puede
reducirse al ingreso, entonces los indicadores monetarios resultan insuficientes para evaluar el
desarrollo. Nussbaum (2000) extiende el enfoque de capacidades de Sen e identifica un conjunto de
capacidades centrales que todo ser humano debe poder ejercer, entre las que destacan la salud
corporal, la integridad física y el control sobre el propio entorno. Desde esta perspectiva, un sistema
de salud deficiente no solo genera costos económicos; compromete la dignidad humana y la
autonomía individual, lo que justifica un tratamiento analítico que va más allá de la eficiencia
asignativa del mercado.

La operacionalización de este enfoque es realizada por Alkire y Foster (2011), quienes construyen
el Índice de Pobreza Multidimensional (IPM) a partir de tres dimensiones: salud, educación y nivel
de vida. La dimensión salud se mide a través de indicadores como la mortalidad infantil y la
desnutrición crónica. Su método identifica como pobre multidimensional a quien carece del umbral
mínimo en al menos un tercio de los indicadores ponderados, lo que permite capturar la intensidad
de la pobreza más allá del ingreso monetario. Esta metodología es también la base del Índice de
Desarrollo Humano (IDH) del Programa de las Naciones Unidas para el Desarrollo, que incorpora
la esperanza de vida al nacer como \textit{proxy} de la dimensión salud. Para sintetizar indicadores
heterogéneos en una medida única, el IDH recurre a técnicas de índices compuestos, entre ellas el
Análisis de Componentes Principales (PCA), que permite reducir la dimensionalidad de un conjunto
de variables correlacionadas preservando la mayor varianza posible (UNDP, 2015).

La relevancia empírica de esta distinción se evidencia con claridad en el caso peruano. El INEI
(2016) reporta que la pobreza monetaria cayó de 58.7\,\% en 2004 a 21.8\,\% en 2015, impulsada
por el crecimiento económico sostenido derivado del boom de materias primas. Sin embargo, como
señala Tezano (2013), la reducción de la pobreza multidimensional fue considerablemente más
lenta, en particular en las regiones de sierra y selva, donde las brechas en acceso a servicios de
salud y saneamiento persistieron incluso durante los años de mayor expansión del PIB. Esta
divergencia sugiere que el crecimiento es condición necesaria pero no suficiente para mejorar el
bienestar en salud. Las diferencias geográficas son marcadas y persistentes: mientras la mortalidad
infantil en Lima se sitúa en niveles próximos a los de países de ingreso medio-alto, en regiones
como Huancavelica y Loreto los indicadores de salud reflejan condiciones propias de economías
mucho más pobres (Ministerio de Salud [MINSA], 2021). Este patrón de heterogeneidad territorial
es uno de los rasgos estructurales más importantes del sector salud peruano y constituye el eje
central del análisis de composición y contexto institucional desarrollado en las secciones posteriores
de este trabajo.

\subsection{La salud como capital humano: efectos sobre la productividad y el crecimiento}

La segunda línea teórica relevante aborda la salud desde la perspectiva del capital humano. La
teoría del capital humano postula que los individuos invierten en educación y salud con el fin de
incrementar su productividad futura (Becker, 1964). Bajo este marco, el gasto en salud no es un
gasto corriente de consumo, sino una inversión que aumenta el stock de capital humano de la
economía y, por tanto, su capacidad productiva de largo plazo.

Grossman (1972) formaliza esta idea en un modelo de demanda de salud en el que los individuos
no demandan atención médica directamente, sino que demandan salud como un stock de capital que
se deprecia con el tiempo y puede ser repuesto mediante inversión: en servicios de salud, nutrición,
ejercicio o conductas preventivas. El modelo predice que la demanda de salud aumenta con el nivel
educativo, porque la educación eleva la eficiencia con que se produce salud, y disminuye con la
edad, porque la tasa de depreciación del capital-salud se eleva conforme el individuo envejece. Este
marco teórico ha sido enormemente influyente en la literatura empírica sobre salud y economía, y
proporciona una base microeconómica sólida para justificar las políticas de prevención y atención
primaria como mecanismos de acumulación de capital humano, distintas en su lógica de las
políticas de atención curativa.

A nivel macroeconómico, esta idea se formaliza en las extensiones del modelo de Solow que
incorporan capital humano como factor de producción: una mejora sostenida en el estado de salud
de la fuerza laboral eleva la productividad total de los factores y desplaza hacia arriba la trayectoria
de crecimiento de estado estacionario (De~Gregorio, 2007). Bloom, Canning y Sevilla (2004)
estiman empíricamente que una mejora de un año en la esperanza de vida de la población adulta
está asociada a un incremento de aproximadamente 4\,\% en el producto per cápita en el largo
plazo, resultado que confirma el papel central de la salud en los modelos de crecimiento endógeno.
Esto implica que el gasto público en salud tiene un efecto potencialmente positivo sobre el
crecimiento del PIB no solo a través del canal de demanda agregada de corto plazo, sino a través de
la acumulación de capital humano en el largo plazo, un efecto que los presupuestos públicos suelen
subestimar al tratar la salud como gasto corriente y no como inversión.

La evidencia para el Perú es analizada empíricamente por Cortez (1999), quien estima el efecto de
las condiciones de salud sobre la productividad laboral e ingresos individuales. Sus resultados
muestran retornos significativos, con diferencias importantes por género y región geográfica: los
efectos son más pronunciados en zonas rurales de sierra y selva, precisamente donde la cobertura de
salud pública es más precaria y la carga de enfermedad es mayor. Estos hallazgos son coherentes
con la literatura internacional sobre trampas de pobreza en salud (Sachs, 2001): una población
enferma es menos productiva, genera menos ingresos, ahorra menos y, en consecuencia, tiene menor
capacidad de invertir en su propia salud, configurando un equilibrio de bajo nivel difícil de romper
sin intervención pública. Este mecanismo de retroalimentación negativa explica en parte la
persistencia de las desigualdades sanitarias regionales en el Perú a pesar del crecimiento
macroeconómico sostenido de la primera década del siglo~XXI.

Valdez et al. (2013) complementan este diagnóstico con el perfil epidemiológico del Perú: las
enfermedades no transmisibles ---cardiovasculares, diabetes, cáncer--- y las enfermedades
diarreicas agudas concentran la mayor carga de morbilidad y constituyen los principales factores
que reducen la productividad laboral. Este hallazgo es consistente con la relación empírica entre
acceso a agua potable segura y mortalidad infantil que se estima en la sección econométrica, dado
que el saneamiento básico es el principal determinante de las enfermedades diarreicas agudas en el
país. La transición epidemiológica incompleta del Perú ---en la que coexisten enfermedades propias
de países pobres con enfermedades propias de países ricos--- representa un desafío específico para
la política sanitaria: el sistema debe atender simultáneamente dos perfiles de demanda con
estructuras de costos y tecnologías de intervención muy distintas, lo que eleva el gasto necesario
para lograr resultados comparables a los de economías en etapas más avanzadas de la transición.

\subsection{Fallas de mercado y el rol del Estado en la provisión de salud}

La justificación teórica de la intervención pública en el sector salud descansa en la presencia de
fallas de mercado que impiden que el equilibrio privado sea eficiente. La primera y más relevante es
la asimetría de información: el paciente no puede evaluar \textit{ex ante} la calidad del servicio
médico ni el precio justo del medicamento, lo que genera una relación de agencia entre el médico y
el paciente y abre espacio para comportamientos oportunistas en la cadena de comercialización
farmacéutica. Arrow (1963), en su artículo fundacional sobre la incertidumbre y la economía del
bienestar médico, señala que esta asimetría de información es tan profunda en el mercado de salud
que el equilibrio competitivo simplemente no puede existir en su forma canónica: el paciente no
puede juzgar la competencia del médico antes ni después del tratamiento, y el médico tiene
incentivos a sobre-prescribir o a recomendar tratamientos más lucrativos. Esta falla justifica la
regulación de precios de medicamentos y la fiscalización por parte de la Dirección General de
Medicamentos, Insumos y Drogas (DIGEMID) en el Perú, cuya función es reducir las rentas
informacionales que capturan las cadenas de farmacias en perjuicio del consumidor final.

A la asimetría de información se añade el problema del riesgo moral en los mercados de seguro de
salud. Cuando el individuo cuenta con cobertura médica, tiende a consumir más servicios de los que
consumiría si pagara el costo total, porque el asegurador asume parte del costo marginal de cada
atención. Esta sobre-utilización eleva el costo promedio del seguro, lo que puede elevar las primas
y excluir a individuos de menor riesgo del mercado. La selección adversa, a su vez, puede conducir
al colapso del mercado de seguros privados en ausencia de intervención pública, tal como lo
formalizan Akerlof (1970) y Rothschild y Stiglitz (1976) en sus modelos de mercados con
información asimétrica. Este es el argumento técnico más sólido para la existencia de un seguro de
salud público y obligatorio: el mercado privado sin regulación generará submercados segmentados
que excluirán sistemáticamente a los más pobres y a los más enfermos, precisamente a quienes más
necesitan cobertura.

La segunda falla son las externalidades positivas: la vacunación y el acceso al saneamiento básico
generan beneficios sociales que exceden ampliamente los beneficios privados, por lo que el mercado
los subprovee en ausencia de intervención pública. Una persona que se vacuna contra la influenza
no solo se protege a sí misma, sino que reduce la probabilidad de transmisión a toda la comunidad:
un beneficio que no se internaliza en la decisión privada de vacunarse. Análogamente, la inversión
en sistemas de agua potable y alcantarillado genera externalidades de salud cuyo valor social supera
con creces el beneficio privado que captura el individuo que accede al servicio. Este argumento
justifica el subsidio público tanto a los programas de inmunización como a la inversión en
saneamiento básico, y explica la correlación negativa entre cobertura de agua potable y mortalidad
infantil que se analiza empíricamente en la sección de series de tiempo.

La tercera falla es el carácter de bien meritorio de la salud: la sociedad considera que el acceso a
servicios sanitarios básicos debe garantizarse con independencia de la capacidad de pago del
individuo, argumento de naturaleza normativa que complementa los anteriores de carácter positivo.
Esta consideración justifica la existencia del Seguro Integral de Salud (SIS) y de EsSalud como
mecanismos de aseguramiento colectivo orientados a garantizar cobertura universal. Cabe señalar
que la fragmentación del sistema de salud peruano ---con cuatro subsistemas principales: SIS,
EsSalud, sanidad de las Fuerzas Armadas y Policía Nacional, y el sector privado--- genera
ineficiencias adicionales derivadas de la duplicación de infraestructura, la ausencia de economías
de escala y la dificultad para la portabilidad de los beneficios entre regímenes (Francke, 2013).
Estas ineficiencias estructurales constituyen, en sí mismas, una falla de coordinación que el diseño
institucional del sector no ha logrado resolver, y que se analiza con mayor detalle en la sección de
contexto institucional.

\subsection{Los determinantes sociales de la salud y la equidad sanitaria}

Un cuarto marco teórico complementario a los anteriores es el enfoque de los determinantes
sociales de la salud, desarrollado por la Comisión sobre Determinantes Sociales de la Salud de la
Organización Mundial de la Salud (OMS, 2008). Este enfoque postula que el estado de salud de la
población no depende únicamente del gasto en servicios médicos ni del comportamiento individual,
sino de un conjunto de condiciones estructurales que incluyen el ingreso, la educación, el empleo,
la vivienda, las redes sociales y el entorno físico. La distribución de la salud en la población no es,
por tanto, un resultado aleatorio, sino el reflejo de la distribución del poder, el ingreso y los
recursos en la sociedad.

Para el Perú, esta perspectiva resulta especialmente relevante dado el patrón de inequidad en la
distribución de la carga de enfermedad. Las poblaciones indígenas, rurales y de menor nivel
educativo presentan tasas de mortalidad infantil, desnutrición crónica y mortalidad materna
significativamente más altas que el promedio nacional, diferencias que no se explican
exclusivamente por el menor gasto en servicios médicos en esas regiones, sino también por la
ausencia de agua potable, saneamiento, infraestructura vial y conectividad (MINSA, 2021).
Wagstaff (2002) proporciona el marco analítico para medir estas inequidades mediante los índices
de concentración, que permiten cuantificar en qué medida la distribución de la salud replica la
distribución del ingreso en una sociedad. Cuando el índice de concentración de la mortalidad
infantil es negativo y significativo ---como ocurre consistentemente en el Perú---, ello indica que
la carga de enfermedad recae desproporcionadamente sobre los más pobres, lo que constituye una
inequidad que el sistema de salud debería reducir como objetivo explícito de política.

La teoría de los determinantes sociales permite, asimismo, evaluar la efectividad de intervenciones
que se ubican fuera del sector salud pero que tienen efectos sanitarios de gran magnitud. Los
programas de transferencias condicionadas como Juntos, que vinculan la asistencia económica a las
familias con el cumplimiento de controles de salud y nutrición, operan como mecanismos de
corrección de los determinantes sociales de la salud sin requerir necesariamente una expansión del
gasto en servicios médicos. De manera análoga, los programas de saneamiento básico rural reducen
la exposición a enfermedades infecciosas actuando sobre el entorno físico, no sobre la capacidad
instalada del sistema sanitario. Esta lógica intersectorial es fundamental para comprender por qué
la mejora de los indicadores de salud en el Perú durante las dos primeras décadas del siglo~XXI
dependió tanto de las políticas macroeconómicas de reducción de la pobreza y del gasto en
infraestructura como del gasto específico en el sector salud.

\subsection{La curva de Preston y la transición epidemiológica peruana}

Un quinto marco teórico que articula el análisis empírico de este trabajo es la curva de Preston
(1975). Preston documentó empíricamente que existe una relación positiva pero cóncava entre el
ingreso per cápita y la esperanza de vida en el corte transversal de países: a medida que el ingreso
crece, la esperanza de vida aumenta, pero a tasas decrecientes. Esta relación tiene dos implicancias
importantes para el caso peruano. En primer lugar, dado que el Perú se ubica en la parte de mayor
pendiente de la curva ---es decir, como país de ingreso medio-bajo, cada dólar adicional de ingreso
per cápita se traduce en una ganancia de esperanza de vida relativamente mayor que en países de
ingreso alto---, existe un espacio considerable para obtener mejoras en los indicadores de salud
mediante inversiones moderadas en tecnología médica y saneamiento. En segundo lugar, la curva de
Preston muestra que su posición se ha desplazado hacia arriba con el tiempo: para un mismo nivel
de ingreso, la esperanza de vida es hoy mayor que hace cincuenta años, lo que refleja el progreso
tecnológico en medicina y salud pública ---vacunas, antibióticos, rehidratación oral--- que ha
reducido la mortalidad independientemente del ingreso.

Este desplazamiento tecnológico ha sido especialmente notable en el Perú: la esperanza de vida al
nacer pasó de aproximadamente 49 años en 1965 a más de 76 años en 2019, un incremento de casi
treinta años en cinco décadas que no puede explicarse únicamente por el crecimiento del ingreso y
que refleja el impacto acumulado de la tecnología médica, la ampliación de la cobertura vacunal y
las mejoras en saneamiento básico (Banco Mundial, 2025). El impacto del COVID-19 en 2020
---que redujo la esperanza de vida en aproximadamente tres años--- ilustra dramáticamente la
magnitud del choque que puede producir una epidemia de gran escala sobre el stock de capital
humano de una economía, y refuerza el argumento a favor de sistemas de salud pública robustos
con capacidad de respuesta ante emergencias sanitarias. El Perú registró una de las tasas de exceso
de mortalidad más altas del mundo durante la pandemia, hecho que pone en evidencia las
fragilidades estructurales de un sistema históricamente subfinanciado y fragmentado.

La transición epidemiológica complementa el análisis de la curva de Preston: conforme los países
reducen la mortalidad por enfermedades infecciosas y mejoran la nutrición, el peso relativo de las
enfermedades crónicas no transmisibles ---cardiovasculares, cáncer, diabetes--- aumenta en la
carga total de morbilidad y mortalidad (Omran, 1971). El Perú se encuentra en una fase intermedia
de esta transición, con una doble carga de enfermedad que exige un sistema de salud capaz de
atender simultáneamente enfermedades propias del subdesarrollo ---tuberculosis, enfermedades
diarreicas, desnutrición--- y enfermedades propias del desarrollo ---hipertensión, diabetes,
cáncer---. Esta dualidad tiene consecuencias directas sobre las necesidades de gasto y la estructura
institucional del sistema de salud. Las enfermedades crónicas no transmisibles requieren atención
continua y de largo plazo, con alta intensidad tecnológica y costo por episodio elevado; las
enfermedades infecciosas y la desnutrición se abordan con mayor eficacia mediante intervenciones
de salud pública preventivas y de bajo costo unitario. Un sistema de salud que no distingue estos
dos perfiles de demanda ni asigna recursos en función de su carga relativa corre el riesgo de ser
simultáneamente ineficiente en el tratamiento de las enfermedades crónicas e insuficiente en la
prevención de las enfermedades infecciosas, resultado que describe razonablemente la situación
actual del sector salud peruano.

% ============================================================
\section*{Referencias}
% ============================================================

\begin{hangparas}{1.5em}{1}

Akerlof, G. (1970). The market for ``lemons'': Quality uncertainty and the market mechanism.
\textit{Quarterly Journal of Economics, 84}(3), 488--500.

Alkire, S., \& Foster, J. (2011). Counting and multidimensional poverty measurement.
\textit{Journal of Public Economics, 95}(7), 476--487.

Arrow, K. (1963). Uncertainty and the welfare economics of medical care.
\textit{American Economic Review, 53}(5), 941--973.

Banco Mundial. (2025). \textit{World Development Indicators}. Washington D.C.: The World Bank.

Becker, G. (1964). \textit{Human Capital: A Theoretical and Empirical Analysis, with Special
Reference to Education}. New York: National Bureau of Economic Research.

Bloom, D., Canning, D., \& Sevilla, J. (2004). The effect of health on economic growth:
A production function approach. \textit{World Development, 32}(1), 1--13.

Cortez, R. (1999). \textit{Salud y productividad en el Perú: Un análisis empírico por género y
región} (Working Paper R-363). Washington D.C.: Inter-American Development Bank.

De Gregorio, J. (2007). \textit{Macroeconomía: Teoría y Políticas}. Santiago: Pearson
Educación.

Francke, P. (2013). \textit{El sistema de salud del Perú: Situación actual y estrategias para
orientar la extensión de la cobertura contributiva}. Lima: OIT.

Grossman, M. (1972). On the concept of health capital and the demand for health.
\textit{Journal of Political Economy, 80}(2), 223--255.

Instituto Nacional de Estadística e Informática. (2016). \textit{Evolución de la Pobreza
Monetaria 2009--2015. Informe Técnico}. Lima: INEI.

Ministerio de Salud. (2021). \textit{Análisis de la Situación de Salud del Perú}. Lima: MINSA.

Nussbaum, M. (2000). \textit{Women and Human Development: The Capabilities Approach}.
Cambridge: Cambridge University Press.

Omran, A. (1971). The epidemiologic transition: A theory of the epidemiology of population
change. \textit{Milbank Memorial Fund Quarterly, 49}(4), 509--538.

Organización Mundial de la Salud. (2008). \textit{Closing the Gap in a Generation: Health
Equity Through Action on the Social Determinants of Health}. Geneva: WHO.

Preston, S. (1975). The changing relation between mortality and level of economic development.
\textit{Population Studies, 29}(2), 231--248.

Rothschild, M., \& Stiglitz, J. (1976). Equilibrium in competitive insurance markets: An essay
on the economics of imperfect information. \textit{Quarterly Journal of Economics, 90}(4),
629--649.

Sachs, J. (2001). \textit{Macroeconomics and Health: Investing in Health for Economic
Development}. Geneva: WHO Commission on Macroeconomics and Health.

Sen, A. (1988). Property and hunger. \textit{Economics and Philosophy, 4}(1), 57--68.

Sen, A. (2000). \textit{Desarrollo y Libertad}. Buenos Aires: Planeta.

Tezano, S. (2013). \textit{Desarrollo humano, pobreza y desigualdades}. Santander: Universidad
de Cantabria.

United Nations Development Programme. (2015). \textit{Human Development Report 2015:
Technical Notes --- Calculating the Indices}. New York: UNDP.

Valdez, W., Napanga, E., Oyola, A., Mariños, J., Vílchez, A., Medina, J., \& Berto, M. (2013).
\textit{Análisis de la Situación de Salud en el Perú}. Lima: Ministerio de Salud.

Wagstaff, A. (2002). \textit{Inequalities in health in developing countries: Swimming against
the tide?} (Policy Research Working Paper No.~2795). Washington D.C.: World Bank.

\end{hangparas}

\end{document}
